\documentclass[11pt,a4paper]{article}
\usepackage[margin=1in]{geometry}
\usepackage{booktabs}
\usepackage{longtable}
\usepackage{graphicx}
\usepackage{hyperref}
\usepackage{listings}
\usepackage{xcolor}
\usepackage{enumitem}
\usepackage{caption}
\usepackage{amsmath}
\usepackage{amssymb}
\usepackage{tcolorbox}
\tcbuselibrary{breakable}

\hypersetup{
    colorlinks=true,
    linkcolor=blue,
    filecolor=magenta,
    urlcolor=cyan,
}

\tcbset{
    promptbox/.style={
        colback=blue!5!white,
        colframe=blue!75!black,
        title={#1},
        fonttitle=\bfseries,
        before skip=6pt,
        after skip=6pt,
        breakable,
    }
}

% Long unbreakable \texttt{} paths are the main source of overfull boxes in
% this report; a generous emergencystretch plus sloppy paragraph breaking lets
% TeX loosen those lines instead of running into the margin.
\setlength{\emergencystretch}{3em}
\sloppy

\title{Performance Report 23: Vulkan Engine \\
\large Texture Arrays, the GPU Mixer, and the Load Path --- the $\sim$930\,MB and $\sim$300 Blocking Submits No Frame Audit Counted}
\author{}
\date{\today}

\begin{document}

\maketitle

\begin{abstract}
Reports~17--22 audited every pass the frame \emph{draws}: ray tracing, the water
machinery and its noise, the solid/water raster core, vegetation, submit counts,
and buffer pools. This report audits what the frame \emph{holds and builds}: the
five material texture arrays, the load path that fills them, the GPU mixer that
regenerates them, and the billboard/editable authoring path around them.

The arrays are $5 \times 32$ layers $\times 1024^2$ RGBA8 with 11 mips each:
\textbf{853\,MiB} of device-local memory and the largest single image commitment
in the engine, which report~22's buffer-oriented tally did not count. 26 of the
32 layers carry data; the other six (160\,MiB) are never written --- five are
referenced by nothing, one is a descriptor-only placeholder. The bring-up path
issues \textbf{$\sim$300 blocking \texttt{runSingleTimeCommands} submits} (each a
fence round-trip polled in 100\,µs sleeps), allocates 105 fresh 4\,MiB staging
buffers that are zero-memset immediately before being overwritten (420\,MiB of
waste), and runs 105 separate blocking mip-generation submits (2{,}310 barriers,
1{,}050 blits). The mixer's every parameter tick --- the UI enqueues on each
slider change --- transitions \emph{all five arrays, all 32 layers, all 11 mips}
into \texttt{GENERAL} and back ($\sim$3{,}500 subresource transitions) and blocks
the CPU on a fence, because the sampler descriptors bind the full array views;
the per-layer descriptor sets written to make generations layer-local are
allocated \emph{before} the per-layer views exist, warn once, and are never
rebuilt (the warning is in every shipped run log). The billboard bake tears
down and recreates all three array images per invocation, and the noise-preview
feature is wired to two mechanisms that are both inert.

It finds \textbf{15 issues} (3 critical, 5 high, 4 medium, 3 low). Each carries
the standard technique and a copy-pasteable prompt preserving validation-clean
execution, Synchronization2, and runtime feature detection. Scope score
\textbf{4.7/10}; the frame-side standing (6.4/10) is unchanged, and folding the
asset slice into the combined estimate at 10\% weight gives \textbf{6.2/10}.
\end{abstract}

\tableofcontents
\newpage

% =====================================================================
\section{Executive Summary}
% =====================================================================

\subsection{Scope and Method}

Four areas, none previously reported on: (1) texture residency --- the five
material arrays sampled by the solid shader, the vegetation arrays, the billboard
arrays, and the impostor capture arrays; (2) the load path
(\texttt{TextureArrayManager::allocate/load/loadTriples}, the staging and
mip-generation helpers it drives); (3) the GPU mixer
(\texttt{services/TextureMixer} + \texttt{shaders/perlin\_noise.comp}) and its
per-generation synchronization; (4) the authoring assets around them
(\texttt{EditableTexture}, \texttt{BillboardService}, \texttt{BillboardCreator},
\texttt{AtlasManager}, \texttt{ImpostorCapture} interfaces). Ray tracing,
water, the raster core, vegetation rendering, frame structure, and the buffer
pools are referenced only at interfaces. Method matches report~22: full reads
with file:line evidence, \texttt{spirv-dis} on shipped modules, byte tallies
from creation code. The application was not executed; dynamic magnitudes are
analytic estimates from verified structure, against the shipped
\texttt{logs/run.log} VRAM counter (3{,}132\,MB device-local at 1080p).

\subsection{Measured Budgets}

\begin{table}[h]
\centering
\small
\begin{tabular}{@{}lrrl@{}}
\toprule
\textbf{Module} & \textbf{Words} & \textbf{Instrs} & \textbf{Role} \\
\midrule
\texttt{perlin\_noise.comp.spv} & 4\,937 & 1\,006 & fBm noise + 5-map blend, one dispatch per generation \\
\texttt{capture.vert.spv} & 697 & 153 & impostor capture VS (60 views, 3 types) \\
\texttt{capture.frag.spv} & 1\,509 & 316 & impostor capture FS (albedo + normal + depth) \\
\bottomrule
\end{tabular}
\caption{Disassembly word counts (\texttt{spirv-dis | wc -w}, the reports~20--22
convention) and instruction counts (op lines). The mixer re-expands fBm and
samples up to ten array taps per texel of a $1024^2$ target.}
\end{table}

Per-bring-up fixed structure: \textbf{$\sim$300 blocking single-time submits},
105 fresh 4\,MiB staging buffers, 105 separate mip submits (2{,}310 barriers,
1{,}050 blits), 66M \texttt{powf} calls for sRGB$\rightarrow$linear, and 135
eager per-layer views/descriptors. Per mixer tick: one blocking submit,
$\sim$135 barrier objects, $\sim$3{,}500 subresource transitions, 50 mip blits.
Committed texture-side memory (\S\ref{sec:vram}): \textbf{$\sim$931\,MiB device +
9\,MiB host}, of which \textbf{160\,MiB is idle layers} and 213\,MiB is the mip
chain itself.

\subsection{Recommended Order}

Texture utilization + content-sized layer count (C1) $\rightarrow$ layer-local
generation and commit-on-release mixer trigger (C2, H4) $\rightarrow$ batched
load path with persistent staging and in-CB mips (C3, H8) $\rightarrow$
in-place billboard/editable updates (H5) $\rightarrow$ LUT / zero-init /
eager-descriptor cleanups (H6, H7) $\rightarrow$ delete the vestigial async
machinery and dead surfaces (M9--M12) $\rightarrow$ hygiene (L13--L15).

% =====================================================================
\section{Cost Model}
% =====================================================================

\begin{equation}
C = \underbrace{\sum_{\mathrm{passes}} G_{\mathrm{pass}} V C_{\mathrm{vert}}}_{\text{\S3, reports 20--21}} + \underbrace{S C_{\mathrm{submit}}}_{\text{report 22}} + \underbrace{L C_{\mathrm{bringup}} + G C_{\mathrm{sweep}} + M_{\mathrm{tex}} C_{\mathrm{pressure}}}_{\text{this report}}
\end{equation}

with $S \approx 17$ submits/frame, $L \approx 300$ blocking submits at bring-up,
$G \approx 1$ generation per UI tick, and $M_{\mathrm{tex}} \approx 0.93$\,GB
resident regardless of scene size. Reports~20--21 attacked the first term and
report~22 the second; this report prices the third, which has two components:
the one-shot bring-up bill (staging churn, blocking submits, CPU decode) and
the recurring residency/sweep bill (what the arrays cost to hold, and what a
single parameter change costs to apply).

% =====================================================================
\section{Texture Residency}
\label{sec:vram}
% =====================================================================

Tallied from creation code at $1920\times1080$ (device memory unless noted).
Mip counts are exact: $1024^2$ arrays carry 11 levels, $512^2$ arrays 10.

\begin{table}[h]
\centering
\small
\begin{tabular}{@{}p{5.9cm}p{4.3cm}r@{}}
\toprule
\textbf{Resource} & \textbf{Sizing} & \textbf{Bytes} \\
\midrule
Main material arrays (albedo/normal/bump/roughness/ao) & $5 \times 32 \times 1024^2$ RGBA8, 11 mips & 853\,MiB dev \\
Vegetation arrays (foliage/grass/wild, 3 maps) & $3 \times 3 \times 512^2$ RGBA8, 10 mips & 12\,MiB dev \\
Billboard arrays (rebuilt per bake) & $3 \times 3 \times 512^2$ RGBA8, 10 mips & 12\,MiB dev \\
Impostor capture arrays & $60 \times 256^2 \times 12$\,B & 45\,MiB dev \\
Composed billboard textures (9 EditableTextures) & $9 \times 512^2$ RGBA8 & 9\,MiB dev + 9\,MiB host \\
\midrule
Total & & $\sim$931\,MiB dev + 9\,MiB host \\
\bottomrule
\end{tabular}
\caption{Committed texture memory. Report~22's buffer-oriented table did not
include the material arrays; they are the largest single image commitment in
the engine. The measured device-local total (3{,}132\,MB in
\texttt{logs/run.log}) is consistent with these tallies plus that report's.}
\end{table}

Of the main arrays' 32 layers, layers 0--20 are loaded from the texture triples,
21--25 are mixer targets (generated at startup), 26 receives views and ImGui
descriptors but no writer, and 27--31 are referenced by nothing: \textbf{six
idle layers, 160\,MiB}. The base levels alone are 640\,MiB; the mip chain adds
another 213\,MiB (33\%).

% =====================================================================
\section{Critical Issues}
% =====================================================================

\subsection{C1: 853\,MiB of Material Arrays, 6 of 32 Layers Idle, No Texture-Side Telemetry}
\label{sec:c1}

\textbf{Severity: Critical.}

The main arrays are sized by a magic constant, not by content:
\texttt{MyApp.cpp:404--406} sets \texttt{layerCount = 32} and calls
\texttt{allocate(layerCount, 1024, 1024)}; \texttt{TextureArrayManager.cpp:202}
derives 11 mips, \texttt{:212} allocates 32 array layers of
\texttt{RGBA8}, \texttt{:216} adds \texttt{STORAGE} usage (for the mixer), and
\texttt{:222} commits it all \texttt{DEVICE\_LOCAL}. Five such arrays are
created (\texttt{:255--260}): $5 \times 32 \times 5.33$\,MiB per layer-map
$= 853$\,MiB. The content is 26 layers ($\sim$693\,MiB); the triples list
(\texttt{MyApp.cpp:408--429}) plus five mixer targets
(\texttt{:442--446}) account for those, and \texttt{MyApp.cpp:466--471} creates
descriptors for a 27th placeholder layer that nothing ever writes. Report~22
C1's \texttt{logMemoryUtilization} (shipped) reports mesh/vertex/index pools,
vegetation chunks, and debug streams --- \emph{not} texture layers or bytes.
The result is the same failure mode report~22 found in buffers, at three times
the size: worst-case capacity, no utilization feedback, no resize path.

\textit{Technique: size from content, log utilization, tier resolution.} The
layer count can be computed from the loaded triples plus the configured mixers
(rounded up with a small paint margin) instead of hard-coded 32; startup should
log used/total layers and bytes with a $<25\%$ idle warning (extending the
shipped report-22 telemetry); a settings tier can halve array resolution
($853 \rightarrow 213$\,MiB) for low-memory presets. Mip levels could be capped
for layers whose max LOD is bounded, but resolution tiering dominates.

\begin{tcolorbox}[promptbox={Prompt --- C1: content-sized texture arrays and utilization telemetry}]
Goal: committed texture memory tracks content. Steps: (1) derive the layer
count from the loaded triples plus configured mixers (+2 paint layers) and pass
it to \texttt{allocate}; (2) extend the startup memory log with per-array
layers used/total and bytes, warning below 25\% utilization; (3) add an array
resolution tier to \texttt{Settings} (default 1024, lower presets 512) flowing
through an idle-guarded reallocate that re-uploads the layers. Files:
\texttt{MyApp.cpp}, \texttt{vulkan/TextureArrayManager.cpp},
\texttt{utils/Settings.hpp}. Validation: identical renders at the default tier,
VRAM delta per tier, validation-clean.
\\ \textbf{Implemented as:} all three steps. (1) The layer count is derived at
bring-up in \texttt{MyApp.cpp}: 21 triples $+$ 5 configured mixers $+$ 2
(editable placeholder $+$ spare) $=$ \textbf{28 layers} instead of the magic
32, with the derivation logged; that is $-107$\,MiB at the default tier
(747 vs 853\,MiB) and the six idle layers shrink to the placeholder plus one
spare. (2) \texttt{TextureArrayManager::logMemoryUtilization} prints
used/allocated layers and bytes (aggregate $+$ per-map breakdown) and warns
below 25\%; it runs at the end of \texttt{setupTextures} and from the scene
renderer's report-22 pool log, so every scene load/generation re-reports the
texture side. (3) \texttt{Settings::textureArraySize} (default 1024; Minimal
512, Maximum 1024) flows through \texttt{MyApp::recreateTextureArrays}: one
device idle (the AGENTS.md major-rebuild case), a
\texttt{TextureArrayManager::recreate} that releases the old per-layer
views/ImGui descriptors, swaps the images and notifies the set-0 listener
exactly once with the new views, a re-upload of the same triples, a mixer
rebind (\texttt{updateComputeDescriptorSets} now writes the same
\texttt{GENERAL} sampler layouts the init path and the pre-dispatch sweep
assume --- its old \texttt{SHADER\_READ\_ONLY} write contradicted them), a
refresh of the five generated layers, and a drain of the deferred destroys
while still idle. Layer indices are resolution-independent, so materials,
mixer targets and the brush/viewer pickers needed no changes; 512$^2$ brings
the five arrays to 187\,MiB. One dependency fix was unavoidable: the tier
routes every source through the resampler, so the 2:1 path is now an
area-average box filter instead of point sampling (also closes M11's
downscale bullet). The tier path is idle-guarded and follows the shipped
water/vegetation rebuild pattern, but it still needs a validation run ---
it touches descriptor sets across the mixer, scene and ImGui in one frame.
\\ \textbf{Postscript (C2):} the mixer rebind step this note mentions was
removed when C2 made the mixer's bindings per-generation ---
\texttt{recreateTextureArrays} now only re-uploads and asks for a fresh
generation; the mixer reads the current views and dimensions at each
generation, so there is nothing persistent to re-point.
\\ \textbf{Postscript (first live run, fixed):} three defects the static
audit could not see. (1) \texttt{TextureArrayManager::allocate} never reset
\texttt{currentLayer}, so a tier rebuild re-uploaded the triples at the old
cursor: the first switch loaded 7 triples into layers 21--27 (over the mixer
targets) and the next found the capacity exhausted and loaded \emph{zero}
(\texttt{1/28} in the utilization log). The cursor is now reset in
\texttt{allocate()}, and \texttt{recreateTextureArrays} logs a mismatch if
the reload cannot place every triple. (2) \texttt{ShadowRenderer}'s
per-cascade sets are built once by copying the shared shadow sets; they were
never refreshed when the arrays were reallocated, so the cascade passes
(\texttt{vkCmdDrawIndexedIndirectCount}) bound the previous, destroyed views
--- VUID-08114 on Set 0, Binding 3 (\texttt{heightArray}). The new
\texttt{ShadowRenderer::refreshCascadeTextureBindings} re-copies bindings
1--18 (minus the per-cascade UBO) from the shared sets, called by
\texttt{updateTextureDescriptorSet} right after it rewrites them. (3) The
static signature now includes the manager's \texttt{version}, so a
reallocation forces the descriptor refresh even if the driver recycles
image/sampler handles, and the update logs any texture/shadow image write it
skips (null view/sampler) instead of letting a stale binding fail at draw
time. The startup telemetry behaved as designed: 27/28 layers (96\%)
initialized on the first run.
\end{tcolorbox}

\subsection{C2: Every Mixer Generation Sweeps All Five Arrays Through GENERAL and Blocks the CPU}
\label{sec:c2}

\textbf{Severity: Critical.}

A texture-mixer parameter change runs the whole pipeline synchronously. The
widget enqueues a generation on \emph{every slider edit}
(\texttt{TextureMixerWidget.cpp:162--178,\allowbreak190--198}); \texttt{MyApp.cpp:4057--4060}
flushes the queue from \texttt{postSubmit()} every frame; and
\texttt{TextureMixer::generatePerlinNoise} submits through
\texttt{runSingleTimeCommands} (\texttt{TextureMixer.cpp:976}), which waits on a
fence polled in 100\,µs sleeps
(\texttt{VulkanApp.cpp:1059--1152,\allowbreak1218--1241}) ---
so a drag stalls the CPU once per frame. The dispatch itself is small
($64\times64$ groups, \texttt{:1098--1101}); the cost is around it:

\begin{itemize}[leftmargin=1.4em,itemsep=2pt]
\item \texttt{:1045--1074}: because the sampler descriptors bind the
\emph{full array views} with \texttt{imageLayout = GENERAL} (the reason is
stated in the comment at \texttt{:1041--1044}), the code transitions every
non-target layer of every generated map to \texttt{GENERAL} before the
dispatch --- a whole-array sweep per array.
\item \texttt{:1197--1228}: the same sweep runs in reverse after, and
\texttt{:1107--1156} regenerates the target layer's five mip chains in the
same command buffer (5 arrays $\times$ 10 blits, 22 barriers each).
\end{itemize}

Per tick: $\sim$135 barrier objects, $\sim$3{,}500 subresource transitions
(32 layers $\times$ 11 mips $\times$ 5 arrays in and out), 50 blits, one
blocking fence round-trip. All 853\,MiB of arrays leave
\texttt{SHADER\_READ\_ONLY} and come back, while the solid pass samples them
every frame.

\textit{Technique: layer-local generation, commit on release.} A generation can
only touch three layers: the write target and the two sampled sources
(\texttt{perlin\_noise.comp:182--200}). Transition exactly those, restore the
two sources, and bind storage through per-layer 2D views so no array-wide
layout change is needed (the machinery for this is already written --- see H4).
Separately, the widget should coalesce edits (e.g.\ generate on
\texttt{IsItemDeactivatedAfterEdit} or at most one generation in flight per
frame) so a slider drag is one sweep, not sixty.

\begin{tcolorbox}[promptbox={Prompt --- C2: layer-local, commit-on-release generation}]
Goal: a parameter change touches three layers, not five arrays. Steps:
(1) rework \texttt{generatePerlinNoise} to transition only target + primary +
secondary subresources and restore the two sources; (2) bind the storage
target's single-layer view (per-layer sets of H4) instead of the array view;
(3) trigger from slider release / one max in-flight generation per frame, not
per edit. Files: \texttt{services/TextureMixer.cpp},
\texttt{widgets/TextureMixerWidget.cpp}, \texttt{shaders/perlin\_noise.comp}.
Validation: identical generated pixels, barrier count per generation down,
validation-clean (watch dynamic-layer-index layout rules; keep the conservative
path behind a feature flag if VVL objects).
\\ \textbf{Implemented as:} all three steps, plus two latent bugs found on the
way. (1) Generation now binds \emph{single-layer} views for everything it
touches --- the target layer as five \texttt{image2D} storage views at
\texttt{GENERAL}, the primary and secondary source layers as ten
\texttt{sampler2D} views at \texttt{SHADER\_READ\_ONLY} (bindings 10--14 are
new; the shader samples every view at layer 0 and stores 2D). One
pre-allocated set is re-pointed per generation (15 writes) and the array-wide
pre/post sweeps are gone: only the target layer's five maps are transitioned
(\texttt{GENERAL} before the dispatch, mip chains after); the sampled layers
never move. Views are re-created on demand, so C1's array rebuild leaves no
stale bindings. (2) The persistent triple/per-map/per-layer set machinery and
\texttt{updateComputeDescriptorSets} were deleted outright --- they were
unreachable and their array-view layout contract was the reason for the sweep
(see the H4/M12 postscripts). (3) The widget commits a parameter edit when
the edit ends: a sticky dirty flag (ImGui reports the slider change during
the drag but not on the release frame, so a bare \texttt{IsAnyItemActive}
gate would miss the final value) plus an inactive-item test; picker clicks
and prev/next switches still generate immediately. Latent bugs fixed: per-map
generation (the preview tabs' \texttt{activeMap}) transitioned one map while
the shader wrote all five --- a layout violation on every UI generation; and
the mixer's cached width/height went stale after C1's resolution tier, so
dispatch size and \texttt{textureSize} are now read live from the manager.
Cost per generation: one blocking submit, 5 layer transitions $+$ 5 mip
chains, no whole-array churn (was $\sim$135 barrier objects and
$\sim$3.5k subresource transitions). The design scopes every descriptor
layout to a single layer view --- the property the old sweep was
compensating for --- but a validation run is still the check that counts.
\end{tcolorbox}

\subsection{C3: The Load Path Runs $\sim$300 Blocking Submits and 420\,MiB of Staging Churn}
\label{sec:c3}

\textbf{Severity: Critical.}

\texttt{TextureArrayManager::load} handles one map at a time
(\texttt{:350--379}): create a fresh 4\,MiB staging buffer
(\texttt{:351}), \texttt{memcpy} the decoded pixels (\texttt{:352}), one
blocking upload submit (\texttt{:354--367}), then a \emph{separate} blocking
mip-generation submit (\texttt{:369--371}). For the 21 triples that is
$21 \times 5 = 105$ uploads and 105 mip submits. Add array creation
(\texttt{allocate}, 5 arrays $\times$ 2 whole-image transitions,
\texttt{:250--251}), the vegetation arrays ($3 \times 5 \times 2 = 30$,
including two constant-color default maps per layer because
\texttt{load()} fills absent paths, \texttt{:382--392}), the five initial mixer
generations, the billboard path (H5: 27 + 12 submits) and the impostor init
(4): \textbf{$\sim$300 blocking submits before the first UI frame}. Every one
of them waits on a fence polled in 100\,µs sleeps
(\texttt{VulkanApp.cpp:1226--1239}), sharing a four-entry command-buffer ring
(\texttt{VulkanApp.hpp:782}) whose slots cannot overlap.

The memory side is no better. \texttt{createBuffer} zero-initializes every
allocation (\texttt{VulkanApp.cpp:4146--4159}); the staging buffer is then
immediately overwritten (\texttt{TextureArrayManager.cpp:352}), so the 105
$\times$ 4\,MiB uploads perform 420\,MiB of zero-memset that no one reads, plus
420\,MiB of allocations/frees that a persistent staging ring would eliminate ---
the exact pattern report~22 M10 shipped for vegetation streaming and the
in-tree \texttt{vulkan/streaming/}\allowbreak\texttt{StagingBufferPool} already
implements.

\textit{Technique: batch per layer, keep staging.} One command buffer per layer
(or per array) recording the five copies and \texttt{recordGenerateMipmaps};
one submit and one fence per layer instead of ten; a persistent staging pool
sized to the peak concurrent upload; \texttt{zeroInit=false} where the caller
overwrites the full range.

\begin{tcolorbox}[promptbox={Prompt --- C3: batch the array bring-up}]
Goal: one or two submits per layer, zero staging churn. Steps: (1) acquire
staging from the existing pool (or a small per-manager pool) instead of
\texttt{createBuffer}/\texttt{destroyBuffer} per map; (2) record all present
maps of a layer into one CB (five \texttt{vkCmdCopyBufferToImage} regions) and
append \texttt{recordGenerateMipmaps} for each map; (3) submit once per layer;
(4) pass \texttt{zeroInit=false} for fully-overwritten uploads. Files:
\texttt{vulkan/TextureArrayManager.cpp}, \texttt{vulkan/VulkanApp.cpp}.
Validation: identical array contents, submit count per load, validation-clean.
\\ \textbf{Implemented as:} all four steps. (1) One persistent host-visible
staging buffer per manager holds all five maps of a layer
($5 \times w \times h \times 4$: 20\,MiB at 1024$^2$, 5\,MiB for the
vegetation arrays); it is acquired once and reused for every layer, so the
bring-up performs zero staging allocations/frees. \texttt{allocate()} and
\texttt{destroy()} release it, so a resolution-tier change re-sizes it on the
next load. (2) \texttt{load()} now decodes every present map straight into
its slice of that buffer (default-filled maps write their constant pattern in
place --- no temporary \texttt{new[]} allocation) and records one command
buffer: five layer transitions, five \texttt{vkCmdCopyBufferToImage} regions,
then \texttt{recordGenerateMipmaps} for each map. (3) One
\texttt{runSingleTimeCommands} per layer: the main bring-up drops from 210
blocking submits (105 copies + 105 mip chains) to \textbf{21}; the vegetation
arrays drop to 3 + defaults. (4) \texttt{createBuffer(..., zeroInit=false)},
so the 420\,MiB of pre-overwrite memset is gone with the churn. This also
lands H8's first step (the mip chains move into the upload CB) --- see the
H8 postscript.
\end{tcolorbox}

% =====================================================================
\section{High-Severity Issues}
% =====================================================================

\subsection{H4: The Per-Layer Descriptor Path That Would Make C2 Cheap Is Dead by Construction}
\label{sec:h4}

\textbf{Severity: High.}

\texttt{TextureMixer::createComputePipeline} ends by allocating one persistent
descriptor set per array layer, each binding that layer's single-layer 2D
storage views, and sets \texttt{hasPerLayerDescSets}
(\texttt{TextureMixer.cpp:444--518}) --- the design that makes a generation
layer-local. It can never succeed on the shipped path: \texttt{MyApp.cpp:437--438}
calls \texttt{init(app, \&textureArrayManager)} (and therefore
\texttt{createComputePipeline}) immediately after \texttt{loadTriples}, while
the per-layer views are created later, at \texttt{MyApp.cpp:448--454,467--471}.
The guard at \texttt{TextureMixer.cpp:456--465} sees empty view vectors, logs
\texttt{Warning: textureArrayManager layer view arrays are smaller than
layerAmount}, and returns. Nothing retries: \texttt{updateComputeDescriptorSets}
(the only other writer of these sets, \texttt{:654--762}) has zero callers in
the tree. The warning is therefore printed on every run --- it is in the shipped
\texttt{logs/run.log:128} and \texttt{logs/gpuav\_run.log:130}.

Every generation consequently falls back to \texttt{tripleComputeDescSet}
(\texttt{:1086--1090}), whose sampler bindings are the full array views at
\texttt{GENERAL} --- which is exactly why C2 must sweep all 32 layers. The fix
for C2 is not just written; it is sitting disabled behind a call-ordering
accident.

\textit{Technique: finish it or delete it.} Re-run the per-layer allocation on
first generation (when the views exist), or create the views before mixer
init; verify that storage writes through single-layer views remove the array
sweep. If the layout rules for dynamically-indexed array sampling force the
conservative sweep under validation, say so in a comment and keep one path.

\begin{tcolorbox}[promptbox={Prompt --- H4: revive the per-layer descriptor sets}]
Goal: the mixer's layer-local path actually runs. Steps: (1) split
\texttt{createComputePipeline} so the per-layer set allocation is a retryable
step; (2) call it on first generation (after \texttt{getImTexture} has created
the views) or create views before mixer init; (3) switch
\texttt{generatePerlinNoise} to \texttt{perLayerDescSets[targetLayer]} and
transition only target/primary/secondary; (4) delete the misleading
empty-vector warning once the failure mode is impossible. Files:
\texttt{services/TextureMixer.cpp}, \texttt{MyApp.cpp}. Validation: generated
pixels identical, warning gone, per-generation barrier count collapses,
validation-clean.
\\ \textbf{Postscript (C2):} resolved by deletion. The persistent per-layer
sets, the unreachable triple/per-map sets and
\texttt{updateComputeDescriptorSets} were removed with the C2 rework;
generation is layer-local by construction (single-layer views, bindings
rewritten per run, views re-created on demand), and the empty-vector warning
went with the code that produced it.
\\ \textbf{Postscript (closure):} verified in-tree --- no per-layer / triple /
per-map machinery or warning remains, and the C2 build's run log no longer
contains it. One last validation-cleanliness gap on this path was fixed with
the closure: \texttt{VulkanApp::recordGenerateMipmaps} records its mip
barriers raw, so only the helper's base-level entry reached the pending
layout map and the tracked state after a generation read
\texttt{TRANSFER\_DST} even though the layer ends in
\texttt{SHADER\_READ\_ONLY}. A second generation of the same target layer
(e.g. two slider commits without a recreate in between) therefore emitted a
pre-barrier whose \texttt{oldLayout} claimed the stale value. The helper now
registers the final per-layer READ state, so repeated generations --- and
C3's batched load, which uses the same helper --- keep the tracker truthful.
H4 is closed.
\end{tcolorbox}

\subsection{H5: Billboard Bake Rebuilds Whole Arrays; Editable Updates Re-Upload Whole Images}
\label{sec:h5}

\textbf{Severity: High.}

\texttt{BillboardService::bakeFromTextures} destroys all three array images and
recreates them on every invocation (\texttt{BillboardService.cpp:80--98,100--137}),
then copies each channel's three layers with one blocking submit \emph{per
layer} (9 submits, \texttt{:145--164}) and generates mips with three more
blocking submits (\texttt{:168--174}). A ``Recompose Billboard'' press costs
nine of those submits immediately (\texttt{BillboardCreator::composeBillboard}
runs three full-image \texttt{updateGPU} calls, \texttt{BillboardCreator.cpp:539--541},
each of which is three blocking single-time passes in
\texttt{EditableTexture.cpp:123--139}: transition, copy, transition); the
startup \texttt{bakeAllBillboards} (\texttt{BillboardCreator.cpp:74--91}) runs
that three times and then the bake --- $\sim$39 blocking submits, three array
teardowns, and a full re-decode of the atlas JPEGs
(\texttt{BillboardCreator.cpp:433--453}). The composed output is $512^2$; the
inputs are entirely unchanged for every layer but the edited one.

\textit{Technique: persistent targets, dirty regions.} Keep the three array
images for the session and copy only changed layers into them in one command
buffer; give \texttt{EditableTexture} a sub-rect upload (the editors know the
dirty bounds); cache decoded atlas pixels across composes.

\begin{tcolorbox}[promptbox={Prompt --- H5: in-place billboard updates}]
Goal: a recompose touches only what changed. Steps: (1) create array images
once (recreate only on size/layer change); (2) record all channel-layer copies
plus mips into one CB and submit once; (3) add dirty-rect tracking to
\texttt{EditableTexture::updateGPU}; (4) cache atlas decodes between composes.
Files: \texttt{services/BillboardService.cpp},
\texttt{widgets/BillboardCreator.cpp}, \texttt{vulkan/EditableTexture.cpp}.
Validation: identical baked arrays, submit count per bake down,
validation-clean.
\\ \textbf{Status (re-iteration, Oct 2026):} still \textbf{OPEN} ---
\texttt{bakeFromTextures} still destroys/recreates all three array images per
call (\texttt{BillboardService.cpp:58--60,91--93}), still issues 9 per-layer
blocking copies plus 3 blocking mip submits; \texttt{EditableTexture::updateGPU}
is still whole-image; \texttt{composeBillboard} still runs $3\times$
\texttt{updateGPU} (\texttt{BillboardCreator.cpp:539--541}) with per-compose
atlas re-decodes and no cache.
\end{tcolorbox}

\subsection{H6: The CPU Load Path Pays 66M \texttt{powf}, 420\,MiB of Dead Memset, and Per-Layer Scans}
\label{sec:h6}

\textbf{Severity: High.}

Three exact, local wastes sit in the load path:

\begin{itemize}[leftmargin=1.4em,itemsep=2pt]
\item \texttt{convertSRGB8ToLinearInPlace} (\texttt{TextureArrayManager.cpp:13--26})
runs three \texttt{std::pow} per texel of every albedo layer: $21 \times
1024^2 \times 3 \approx 66$M transcendental calls. Input and output are both
8-bit, so the mapping has 256 possible results --- a 256-entry LUT gives
bit-identical output (same formula, same float rounding, precomputed once) at
$\sim$1/60th the cost. (The call is invoked only for the albedo map,
\texttt{:408--410}.)
\item \texttt{createBuffer} zero-memsets host-visible buffers before returning
(\texttt{VulkanApp.cpp:4146--4159}); \texttt{load()} overwrites the full staging
range immediately after (\texttt{:351--352}). 105 uploads $\times$ 4\,MiB
$=$ 420\,MiB of memset whose only readers are the allocator.
\item \texttt{meanLinearRGB} (\texttt{:305--320}) double-accumulates over
$1024^2$ texels for every albedo layer (21 scans) to maintain a 3-float
average used by the RT proxy.
\end{itemize}

\textit{Technique: LUT, \texttt{zeroInit=false}, fold the scan.} One 256-byte
table replaces the pow loop exactly; the zero-init flag already exists; the
average can be accumulated in the same pass that converts (or from a 256-bin
histogram, since the data is quantized).

\begin{tcolorbox}[promptbox={Prompt --- H6: exact, local load-path waste}]
Goal: remove work with a proof, not a benchmark. Steps: (1) replace
\texttt{convertSRGB8ToLinearInPlace}'s pow loop with a static 256-entry LUT
computed by the same expression; (2) call \texttt{createBuffer(...,
zeroInit=false)} in the upload helper (the memcpy covers the whole range);
(3) accumulate the albedo average while converting. Files:
\texttt{vulkan/TextureArrayManager.cpp}. Validation: pixel-identical arrays
(LUT is exact), CPU time in load, validation-clean.
\\ \textbf{Status (re-iteration, Oct 2026):} \textbf{PARTIAL} --- the pow loop
is still per-texel (\texttt{TextureArrayManager.cpp:15--28}, no LUT) and the
mean is still a separate full scan; only the memset third is fixed, and only
on the batched main/vegetation load via C3's persistent staging
(\texttt{zeroInit=false}). Other uploaders (e.g.\ \texttt{EditableTexture})
still take the default zero-init.
\end{tcolorbox}

\subsection{H7: 135 Per-Layer Views and ImGui Descriptors Are Created Eagerly}
\label{sec:h7}

\textbf{Severity: High.}

\texttt{TextureArrayManager::getImTexture} is explicitly lazy --- it creates the
2D view and \texttt{ImGui\_ImplVulkan\_AddTexture} descriptor on first request
(\texttt{TextureArrayManager.cpp:749--802}). \texttt{MyApp.cpp:448--454} defeats
that by calling it for \emph{all} 26 populated layers $\times$ 5 maps at
startup, and \texttt{:466--471} adds the placeholder layer: \textbf{135 views
and 135 ImGui descriptor sets} before the first frame, most of which the UI
never displays. Each is tracked in \texttt{VulkanResourceManager}, invalidated
and recreated across swapchain recreation (\texttt{MyApp.cpp:3318--3337}), and
deferred-destroyed at shutdown. The TextureViewer/Mixer/Brush3d widgets already
call \texttt{getImTexture} on demand, so the eager loop buys nothing.

\textit{Technique: delete the warm-up loop.} Keep lazy creation; create at most
the visible previews.

\begin{tcolorbox}[promptbox={Prompt --- H7: stop warming every preview}]
Goal: views and descriptors exist only if shown. Steps: remove the eager loops
in \texttt{MyApp.cpp:448--454} and \texttt{:466--471} (keep the layer
initialized flags where the UI needs them); confirm the widgets' on-demand path
covers the first display frame. Files: \texttt{MyApp.cpp}. Validation:
descriptor count at startup, previews still appear, validation-clean.
\\ \textbf{Status (re-iteration, Oct 2026):} still \textbf{OPEN} --- the eager
loops are still present (\texttt{MyApp.cpp:475--481} all loaded layers
$\times$ 5 maps, plus the placeholder at \texttt{:492--497}).
\end{tcolorbox}

\subsection{H8: Mipmaps Are One Blocking Submit per Map --- 105 of Them}
\label{sec:h8}

\textbf{Severity: High.}

\texttt{load()} calls the blocking \texttt{generateMipmaps} per map
(\texttt{TextureArrayManager.cpp:369--371}); each call re-queries format
properties and opens a fresh single-time submit
(\texttt{VulkanApp.cpp:1339--1352}). Inside,
\texttt{recordGenerateMipmaps} (\texttt{:1367--1488}) records, per layer, 22
barriers and 10 blits for 11 levels. At bring-up that is 105 submits,
\textbf{2\,310 barriers and 1\,050 blits} for the main arrays (plus 285/135 for
vegetation). The mixer already shows the better shape: it records
\texttt{recordGenerateMipmaps} \emph{inside} its generation command buffer
(\texttt{TextureMixer.cpp:1142--1155}), one submit for dispatch + five mip
chains. The same is available to the load path the moment it batches (C3);
additionally the format-capability query can be cached once instead of 105
times.

\textit{Technique: mips in the upload command buffer.}

\begin{tcolorbox}[promptbox={Prompt --- H8: fold mip generation into the upload submit}]
Goal: no separate mip submits at load. Steps: (1) after the five copies, call
\texttt{recordGenerateMipmaps} in the same CB per map (C3's batching makes this
free); (2) cache the linear-blit format check per format; (3) update tracked
layouts once per layer. Files: \texttt{vulkan/TextureArrayManager.cpp},
\texttt{vulkan/VulkanApp.cpp}. Validation: identical mip content, submit count
per layer from 10 to 1, validation-clean.
\\ \textbf{Postscript (C3):} shipped with the C3 batch --- the per-map mip
chains are recorded inside the per-layer upload command buffer, so the 105
separate blocking mip submits are gone, and with them the per-call
format-capability queries of the blocking \texttt{generateMipmaps} wrapper
(the load path no longer calls it).
\end{tcolorbox}

% =====================================================================
\section{Medium-Severity Issues}
% =====================================================================

\subsection{M9: Vestigial Async Generation Machinery Runs Every Frame}
\label{sec:m9}

\textbf{Severity: Medium.}

\texttt{TextureMixer} tracks ``pending fences'' for asynchronous generations,
but no code ever inserts into \texttt{pendingFences} --- every reference in the
tree reads or clears it (\texttt{TextureMixer.cpp:128--161,217--253}). The
consequences: \texttt{isLayerGenerationPending} always returns false,
\texttt{waitForLayerGeneration} is unreachable in practice,
\texttt{getPendingGenerationCount} counts only queued requests, and
\texttt{pollPendingGenerations} iterates an always-empty vector. The last one
still runs \emph{every frame} (\texttt{MyApp.cpp:4059--4060}) and calls
\texttt{app->processPendingCommandBuffers()} (\texttt{TextureMixer.cpp:165}),
which \texttt{drawFrame} already called immediately before
\texttt{postSubmit()} (\texttt{VulkanApp.cpp:5044}) --- a duplicate walk of the
deferred-destroy lists per frame. Generation has been synchronous since the
\texttt{runSingleTimeCommands} rewrite; the async scaffolding predates it.

\textit{Technique: delete or complete.} Either remove the fence machinery and
its per-frame calls, or make generation genuinely async (the fences would then
be populated by the submit), which would also unblock C2's commit-on-release
design.

\begin{tcolorbox}[promptbox={Prompt --- M9: one story for mixer synchronization}]
Goal: no dead synchronization paths. Steps: (1) decide synchronous or async;
if synchronous, delete \texttt{pendingFences}, the wait/poll/is-pending
accessors, and the \texttt{processPendingCommandBuffers} call in
\texttt{pollPendingGenerations}; (2) keep the per-frame flush;
(3) update the widget's pending counter to the queue depth. Files:
\texttt{services/TextureMixer.cpp,.hpp}, \texttt{MyApp.cpp}. Validation:
identical behavior, no duplicate pending-CB walk, validation-clean.
\\ \textbf{Status (re-iteration, Oct 2026):} still \textbf{OPEN} --- the fence
machinery is retained (\texttt{TextureMixer.cpp:128--143,191,205--206,218--230}),
still polled every frame from \texttt{postSubmit} (\texttt{MyApp.cpp:4148--4152}),
still never populated.
\end{tcolorbox}

\subsection{M10: Five Initial Generations Are Five Full Sweeps}
\label{sec:m10}

\textbf{Severity: Medium.}

\texttt{generateInitialTextures} enqueues five mixers
(\texttt{MyApp.cpp:442--446,474}); the first \texttt{postSubmit} flushes them
sequentially, and each one performs C2's whole-array sweep, its own blocking
submit, and its own five mip chains --- five identical layout round-trips over
the same five arrays in one frame. The dispatches are independent (different
target layers, disjoint storage), so they can share one command buffer, one
sweep, and one fence.

\textit{Technique: coalesce same-frame generations.} Flush all pending
requests into one command buffer; transition each affected array once
(union of targets and sources); dispatch per request.

\begin{tcolorbox}[promptbox={Prompt --- M10: batch the initial generations}]
Goal: N queued generations cost one sweep and one submit. Steps: build a union
of target/source layers per array, emit the transitions once, record one
dispatch per request (push constants differ), append each target layer's mip
chain, submit once. Files: \texttt{services/TextureMixer.cpp}. Validation:
identical generated layers, startup generation cost down, validation-clean.
\\ \textbf{Status (re-iteration, Oct 2026):} still \textbf{OPEN} ---
\texttt{generateInitialTextures} still enqueues sequentially
(\texttt{TextureMixer.cpp:58--78}) and the flush still runs one blocking
submit plus five mip chains per generation. C2's layer-local transitions make
the union-sweep variant cheaper than described, but it is not implemented.
\end{tcolorbox}

\subsection{M11: Vegetation Sources Are Nearest-Downscaled; Mixer Layers Keep Gray RT Albedo}
\label{sec:m11}

\textbf{Severity: Medium.}

Two content-path gaps with small but real cost/quality effects:

\begin{itemize}[leftmargin=1.4em,itemsep=2pt]
\item The vegetation arrays are $512^2$ (\texttt{MyApp.cpp:3496}) while all
three atlas sources are $1024^2$; \texttt{load()} therefore routes them
through \texttt{resizeNearest} (\texttt{TextureArrayManager.cpp:400--406}), a
point-sampled 2:1 decimation of alpha-cut foliage --- the classic aliasing
recipe for wind-animated billboards. A box filter or mip-filtered downscale
costs the same at load and is strictly better at render time.
\item \texttt{albedoAvg} (the RT proxy's per-layer albedo) is only written by
\texttt{load()} (\texttt{:416--417}) and the dead
\texttt{updateLayerFromEditableMap} (\texttt{:707--720}); mixer-generated
layers 21--25 keep the gray $(0.5,0.5,0.5)$ default
(\texttt{:281}), so reflection/refraction rays that hit mixed terrain return
gray proxies.
\end{itemize}

\textit{Technique: filtered downscale; update the average on generation.} Use a
2$\times$2 box (or \texttt{stbir\_resize}) for resampling; compute the average
in the generation path from the same data the compute shader blended (or read
back one mip).

\begin{tcolorbox}[promptbox={Prompt --- M11: filtered vegetation downscale, live RT albedo}]
Goal: no point-decimation, no stale proxy albedo. Steps: (1) replace
\texttt{resizeNearest} with a box/bilinear resample (or load at target size via
stb's resize); (2) update \texttt{albedoAvg} for mixer target layers after
generation. Files: \texttt{vulkan/TextureArrayManager.cpp},
\texttt{services/TextureMixer.cpp}. Validation: vegetation A/B (aliasing),
reflection color spot-check, validation-clean.
\\ \textbf{Status (re-iteration, Oct 2026):} \textbf{PARTIAL} --- the downscale
half is fixed (box-filter \texttt{resizeAreaAverage},
\texttt{TextureArrayManager.cpp:404,494}; \texttt{resizeNearest} is gone from
the tree). The albedo half is open: the generation path still never writes
\texttt{albedoAvg}, so mixer layers keep the gray default
(\texttt{TextureArrayManager.cpp:356--357}).
\end{tcolorbox}

\subsection{M12: Dead Mixer Surfaces and Duplicated Samplers}
\label{sec:m12}

\textbf{Severity: Medium.}

The mixer carries several write-only surfaces: the five per-map descriptor sets
allocated in \texttt{createComputeDescriptorSet} (\texttt{TextureMixer.cpp:438--442,804--879})
are never bound (\texttt{generatePerlinNoise} always uses the triple set);
\texttt{updateComputeDescriptorSets} has no callers and binds samplers at
\texttt{SHADER\_READ\_ONLY} while the init-time set uses \texttt{GENERAL} ---
two contradictory descriptions of the same bindings; and
\texttt{TextureArrayManager::updateLayerFromEditableMap} (56 lines of
transition/copy/mip choreography plus the only caller of
\texttt{waitForLayerGeneration}) has zero callers. Ten samplers are created
with identical state (five per manager, \texttt{VulkanApp::createTextureSampler}
at \texttt{TextureArrayManager.cpp:269--273} and \texttt{MyApp.cpp:3496}), and
\texttt{TextureArrayManager::create()} (\texttt{:488--542}) is a whole
zero-layer initializer with no callers. Each item is small; together they are
the reason C2/H4's intended path is hard to see.

\textit{Technique: delete the unbound paths, share the samplers.} One sampler
per configuration, one descriptor-set recipe.

\begin{tcolorbox}[promptbox={Prompt --- M12: one recipe per purpose}]
Goal: every surface has a caller. Steps: (1) delete the per-map sets and
\texttt{updateComputeDescriptorSets} (or make them the one path and delete the
triple fallback); (2) delete \texttt{updateLayerFromEditable(Map)} and
\texttt{create()} if the tree still shows no callers; (3) create each sampler
kind once and share. Files: \texttt{services/TextureMixer.cpp,.hpp},
\texttt{vulkan/TextureArrayManager.cpp,.hpp}. Validation: build minus the
symbols, identical behavior, validation-clean.
\\ \textbf{Postscript (C2):} the mixer-side surfaces (triple / per-map /
per-layer descriptor sets, the per-map creation helper and
\texttt{updateComputeDescriptorSets}) were deleted with the C2 rework;
\texttt{TextureArrayManager::create()} and the
\texttt{updateLayerFromEditable(Map)} pair remain as listed.
\\ \textbf{Status (re-iteration, Oct 2026):} unchanged --- \texttt{create()}
(\texttt{TextureArrayManager.cpp:616--669}) and
\texttt{updateLayerFromEditable(Map)} (\texttt{:672--752}) still have zero
callers, and the per-manager samplers are still duplicated
(\texttt{TextureArrayManager.cpp:345--349} plus
\texttt{BillboardService.cpp:176--193}).
\end{tcolorbox}

% =====================================================================
\section{Low-Severity / Hygiene Issues}
% =====================================================================

\subsection{L13: The Noise Preview Is Broken Twice Over}
\label{sec:l13}

\texttt{TextureMixer::getNoiseDescriptor} renders the mask by binding an
alpha-swizzled view (\texttt{TextureMixer.cpp:789--802}), but
\texttt{TextureArrayManager::getImTextureAlpha} returns the cached
\texttt{albedoImTextures[layer]} if it exists
(\texttt{TextureArrayManager.cpp:831}) --- and H7 guarantees it exists, because
startup warmed every layer. The swizzle is never created. Independently, the
shader's \texttt{debugOutput} push constant is never assigned by the C++
(\texttt{PerlinPushConstants.hpp:19} is left to an uninitialized local in
\texttt{TextureMixer.cpp:916}) and is never read by the shader body
(\texttt{shaders/perlin\_noise.comp:202--215} always writes the blend), so the
``Show noise'' checkbox does nothing on either path. Harmless to performance;
worth fixing while touching the mixer.

\textbf{Status (re-iteration, Oct 2026):} still \textbf{OPEN} --- the
early-return still precedes the swizzle
(\texttt{TextureArrayManager.cpp:936--959}) and \texttt{debugOutput} is still
unassigned and unread (\texttt{PerlinPushConstants.hpp:9--10,21},
\texttt{shaders/perlin\_noise.comp:43}).

\subsection{L14: Per-Generation Log String Churn}
\label{sec:l14}

Every enqueue, replacement, and immediate generation formats and appends a
string under \texttt{logsMutex} (\texttt{TextureMixer.cpp:95--100,180--185,883--888}),
and the widget polls the queue depth with two mutex acquisitions per frame
(\texttt{:203--207}). During a slider drag this is one formatted string per
frame for a message no one reads. Gate on a log flag or drop the per-request
lines.

\textbf{Status (re-iteration, Oct 2026):} still \textbf{OPEN} --- the
\texttt{snprintf}+\texttt{emplace} under \texttt{logsMutex} per
enqueue/replace/immediate generation is unchanged
(\texttt{TextureMixer.cpp:96--99} and neighbors); C2's commit-on-release only
lowered how often it fires.

\subsection{L15: Magic 32, Stale Comments, and Misleading Startup Text}
\label{sec:l15}

\texttt{MyApp.cpp:404}'s \texttt{layerCount = 32} is unrelated to the 26 layers
actually used (C1); the mixer's startup banner says ``Albedo, Normal, Bump''
while five maps are generated (\texttt{TextureMixer.cpp:65}); the comment above
\texttt{getImTextureAlpha} promises swizzled alpha that never happens (L13);
and \texttt{TextureArrayManager.hpp:83}'s ``version counter'' is bumped but
nothing in the tree reads it. Update the comments, delete the counter if
unused, and let C1's content-derived count replace the constant.

\textbf{Status (re-iteration, Oct 2026):} \textbf{PARTIAL} --- the magic 32 is
gone (derived count, \texttt{MyApp.cpp:452}) and the version counter is now
read (\texttt{SceneRenderer.cpp:1140}, shadow cascade refresh). Still open:
the ``Albedo, Normal, Bump'' banner for five maps
(\texttt{TextureMixer.cpp:65}) and the alpha-swizzle comment L13 never
delivers (\texttt{TextureArrayManager.hpp:127--129}).

% =====================================================================
\section{Implementation Roadmap and Expected Deltas}
% =====================================================================

\begin{table}[h]
\centering
\small
\begin{tabular}{@{}lp{5.8cm}p{5.6cm}@{}}
\toprule
\textbf{Item} & \textbf{Change} & \textbf{Est.\ effect} \\
\midrule
C1 & Content-sized layers + utilization log + resolution tier & $-160$\,MiB idle; $-640$\,MiB at the 512 tier \\
C2 & Layer-local, commit-on-release generation & $-3.5$k subresource transitions/tick, no per-tick fence stall \\
C3 & Batched load path, persistent staging, \texttt{zeroInit=false} & $-210$ submits, $-420$\,MiB memset, $-420$\,MiB alloc churn \\
H4 & Revive per-layer descriptor sets & prerequisite for C2 \\
H5 & In-place billboard/editable updates & $\sim$39 $\rightarrow$ $\sim$3 submits per bake, no array rebuild \\
H6 & sRGB LUT + no dead memset + fused average & $-66$M \texttt{powf}, $-420$\,MiB memset \\
H7 & Drop the eager preview loop & $-135$ views/descriptors at startup \\
H8 & Mips in the upload CB & $-105$ submits, $-2.3$k barriers \\
M9--M12 & One sync story, coalesced generations, filtered downscale, dead-surface sweep & additive \\
L13--L15 & Noise preview, log churn, constants/comments & noise floor \\
\bottomrule
\end{tabular}
\end{table}

C1--C3 are the items that change what the engine \emph{costs to run} on the
texture side: residency, interactive stall, and bring-up. H4 is not optional
for C2 --- it was the intended mechanism. None of the changes touches rendered
content except M11's downscale filter (a quality fix) and none changes the
lighting model.

% =====================================================================
\section{Overall Score}
\label{sec:score23}
% =====================================================================

\subsection{Rubric (Report-23 Scope)}

\begin{table}[h]
\centering
\small
\begin{tabular}{@{}lcp{7.2cm}@{}}
\toprule
\textbf{Category} & \textbf{Wt.} & \textbf{Score and justification} \\
\midrule
Texture residency & 30\% & \textbf{4.5} --- arrays are the right structure
(mips, aniso, deferred frees, few descriptors), but 853\,MiB is provisioned by
a magic constant, 160\,MiB is idle, and no texture-side telemetry exists. \\
Generation path & 25\% & \textbf{4.0} --- one blocking fence round-trip and a
five-array \texttt{GENERAL} sweep per slider tick; the layer-local fix is
written and disabled by call ordering. \\
Load / bring-up path & 25\% & \textbf{5.0} --- correct and one-shot, but
$\sim$300 blocking submits, per-map staging churn, per-map mip submits, and
420\,MiB of dead memset. \\
Authoring assets & 20\% & \textbf{5.5} --- atlas autodetect, nearest-leaf fill
and lazy previews are good; bakes rebuild whole arrays, edits re-upload whole
images, previews are eager. \\
\bottomrule
\end{tabular}
\end{table}

\subsection{Overall: 4.7/10 Scope, 6.2/10 Combined}

\begin{equation}
0.30(4.5) + 0.25(4.0) + 0.25(5.0) + 0.20(5.5) = \mathbf{4.7}
\end{equation}

Report~22 scored the frame-side scope 6.0 and the combined engine 6.4; that
estimate implicitly valued the asset/load dimension at roughly 10\% without
measuring it. Measuring it:

\begin{equation}
0.90(6.4) + 0.10(4.7) = \mathbf{6.2}
\end{equation}

The engine's verdict is unchanged in kind: the architecture (arrays, bindless
sizing, explicit layout tracking, deferred destruction, one-shot compute
generation) is sound; the constant factors --- capacity, synchronization
granularity, and batching --- are again where the money is. Landing C1--C3 and
H4 moves this scope toward $\sim$8/10 with no redesign.

\end{document}
